Báo cáo này khảo sát bài toán phân đoạn nguyên liệu thực phẩm fine-grained trên bộ dữ liệu \texttt{FoodSeg103} dưới góc nhìn lightweight segmentation, xuất phát từ yêu cầu cân bằng giữa chất lượng phân đoạn và hiệu quả suy luận trong các kịch bản triển khai tài nguyên hạn chế. Phân tích dữ liệu tiền huấn luyện cho thấy \texttt{FoodSeg103} có phân bố long-tail mạnh, ranh giới mờ, chồng lấp cao và tương đồng liên lớp lớn, khiến mô hình nhẹ dễ thiên lệch và suy giảm chi tiết cục bộ.

Để thiết lập mốc tham chiếu, BiSeNet được huấn luyện \emph{trực tiếp} trên không gian nhãn nguyên thủy gồm 103 lớp của \texttt{FoodSeg103} mà không áp dụng bất kỳ bước tái cấu trúc nhãn nào. Pipeline huấn luyện, augmentation và đánh giá nhất quán theo mIoU, mAcc và pixel accuracy nhằm làm rõ mức độ phù hợp của kiến trúc lightweight đối với bài toán này. Kết quả baseline ($\mathrm{mIoU}=0.31$) cung cấp mốc tham chiếu rõ ràng cho các cải tiến texture-aware và boundary-aware ở giai đoạn tiếp theo.

Sau khi phân tích lỗi baseline, một thực nghiệm mở rộng độc lập được thực hiện: tái cấu trúc không gian nhãn từ 103 xuống 75 lớp dựa trên tương đồng texture GLCM và quan hệ ngữ nghĩa, tạo ra bộ \texttt{FoodSeg103\_rebalanced}. Thực nghiệm này đóng vai trò phân tích bổ sung để làm rõ ảnh hưởng của phân mảnh nhãn đến lỗi của mô hình, không thay thế đánh giá trên bài toán gốc và cũng không được dùng làm dữ liệu huấn luyện cho baseline.

Mã nguồn thực nghiệm được công bố tại \url{https://github.com/ariesanhthu/SemanticSegmentation-FoodSeg103}.